\section{BT1089--BT1090 Natural Core and Reservoir Intertwiner}
\label{sec:bt1089-bt1090-natural-core-intertwiner}

The arbitrary ordered-basis choice in the first $40\hookrightarrow54$ skeleton can be removed.  Identify the right-chirality matter sheet in one generation with the regular matter-shell address cube
\[
C^{27}_R(g)=\mathbb C[\mathbb F_3^3].
\]
After the parabolic shell origin is fixed, the antipodal involution $x\mapsto -x$ has one fixed point and thirteen two-point orbits on the nonzero vectors.  Define
\[
B^{13}_R(g)=\{f\in\mathbb C[\mathbb F_3^3]: f(-x)=-f(x)\}.
\]
Equivalently, $B^{13}_R(g)$ is spanned by $e_x-e_{-x}$ over the thirteen projective directions in $PG(2,3)$.  The residual is the even subspace
\[
R_{14}(g)=\{f:f(-x)=f(x)\},
\]
spanned by the origin together with the thirteen pair-sums.  Thus
\[
\dim B^{13}_R(g)=13,
\qquad
\dim R_{14}(g)=14,
\qquad
C^{40}_{\mathrm{core}}(g)=C^{27}_L(g)\oplus B^{13}_R(g).
\]
The $13=\Phi_3$ is therefore no longer a coordinate convention: it is the number of projective directions in the ternary matter cube.

The reservoir coupling also has its first nonzero candidate.  Each lifted $P_{22}(g)$ splits as
\[
P_{22}(g)=F_{13}(g)\oplus D_9(g),
\]
where $F_{13}$ is the fixed/gauge-perp part and $D_9$ is the shell-cycle-sum diagonal matter part.  The fixed block has $13=1+12$; remove its trace line and identify the trace-free quotient with the gauge packet layer
\[
A_{12}=A_1\oplus A_3\oplus A_8.
\]
The first reservoir intertwiner is
\[
K:T_{66}\to A_{12},
\qquad
K=\frac13\sum_{g=0}^2 \pi_{12}^{(g)}\operatorname{pr}_{F_{13}(g)}.
\]
It has rank $12$, while its adjoint embeds a gauge packet diagonally back into the three fixed/gauge-perp blocks.  The $D_9(g)$ blocks remain diagonal matter bookkeeping at this stage.
